% !TEX root = ../../../main.tex
% 所属：短篇 / 短篇1：Weierstrass第一逼近定理的几种证明
% 内容：正文（原文未分节）
% 正文按原文件顺序保存；由 ../article1.tex 调用。
\begin{definition}(Bernstein 多项式)
	$Y$ 是多项式全体构成的集合,定义映射
	\[
	B_n:C[0,1]\longrightarrow Y
	\]
	为
	\[
	f(t)\longmapsto B_n(f,x)=\sum_{k=0}^{n} f\left(\frac{k}{n}\right) C_n^k x^k(1-x)^{n-k}.
	\]
	这里 $B_n(f,x)$ 表示 $f\in X$ 在映射 $B_n$ 作用下的像,它是以 $x$ 为变量的 $n$ 次多项式,称为 $f$ 的 Bernstein 多项式.
\end{definition}

关于映射 $B_n$,直接从定义出发,可以证明它有下述基本性质与基本关系.

对 $B_n$ 有
\begin{enumerate}
	\item $B_n$ 是线性映射,即对任意 $f,g\in X$ 与 $\alpha,\beta\in\mathbb{R}$,成立
	\[
	B_n(\alpha f+\beta g,x)=\alpha B_n(f,x)+\beta B_n(g,x).
	\]
	\item $B_n$ 具有单调性,即若 $f(t)\geqslant g(t)$ 对一切 $t\in[0,1]$ 成立,则
	\[
	B_n(f,x)\geqslant B_n(g,x)
	\]
	对一切 $x\in[0,1]$ 成立.
	\item 对一切 $x\in[0,1]$,成立
	\[
	B_n(1,x)=1,\qquad B_n(t,x)=x,\qquad B_n(t^2,x)=x^2+\frac{x-x^2}{n}.
	\]
	进一步,对任意常数 $s$,有
	\[
	B_n((t-s)^2,x)=\frac{x-x^2}{n}+(x-s)^2.
	\]
\end{enumerate}



\begin{theorem}(Weierstrass 第一逼近定理)\label{weierstrass第一逼近定理}
	设 $f(x)$ 是闭区间 $[a,b]$ 上的连续函数,则对任意给定的 $\varepsilon>0$,存在多项式 $P(x)$,使得
	\[
	|P(x)-f(x)|<\varepsilon
	\]
	对一切 $x\in[a,b]$ 成立.
\end{theorem}
这个证明来自于苏联的Korovkin
\begin{proof}
	一般性地,我们先证 $[a,b]=[0,1]$ 的情形.
	
	由于 $f$ 在 $[0,1]$ 上连续,所以必定有界,即存在 $M>0$,使得对一切 $t\in[0,1]$ 成立
	\[
	|f(t)|\leqslant M.
	\]
	又由 Cantor 定理, $f$ 在 $[0,1]$ 上一致连续.于是对任意给定的 $\varepsilon>0$,存在 $\delta>0$,使得对一切 $t,s\in[0,1]$,当 $|t-s|<\delta$ 时,成立
	\[
	|f(t)-f(s)|<\frac{\varepsilon}{2}.
	\]
	当 $|t-s|\geqslant\delta$ 时,有
	\[
	|f(t)-f(s)|\leqslant 2M\leqslant\frac{2M}{\delta^2}(t-s)^2.
	\]
	也就是说,对一切 $t,s\in[0,1]$,成立
	\[
	-\frac{\varepsilon}{2}-\frac{2M}{\delta^2}(t-s)^2
	\leqslant f(t)-f(s)\leqslant
	\frac{\varepsilon}{2}+\frac{2M}{\delta^2}(t-s)^2.
	\]
	把上式看作关于 $t$ 的连续函数不等式,并对其作用 $B_n$,由 $B_n$ 的单调性得到对一切 $x,s\in[0,1]$,
	\[
	-\frac{\varepsilon}{2}-\frac{2M}{\delta^2}\left[\frac{x-x^2}{n}+(x-s)^2\right]
	\leqslant B_n(f,x)-f(s)
	\leqslant
	\frac{\varepsilon}{2}+\frac{2M}{\delta^2}\left[\frac{x-x^2}{n}+(x-s)^2\right].
	\]
	令 $s=x$,并注意到 $x(1-x)\leqslant\dfrac14$,即得
	\[
	\left|\sum_{k=0}^{n} f\left(\frac{k}{n}\right)C_n^k x^k(1-x)^{n-k}-f(x)\right|
	\leqslant \frac{\varepsilon}{2}+\frac{M}{2n\delta^2}.
	\]
	取
	\[
	N=\left[\frac{M}{\delta^2\varepsilon}\right]+1,
	\]
	当 $n>N$ 时,有
	\[
	\left|\sum_{k=0}^{n} f\left(\frac{k}{n}\right)C_n^k x^k(1-x)^{n-k}-f(x)\right|<\varepsilon
	\]
	对一切 $x\in[0,1]$ 成立.而左边的求和正是多项式 $B_n(f,x)$ 与 $f(x)$ 的差.
	
	一般闭区间 $[a,b]$ 的情形,只需对区间作线性变换化为 $[0,1]$ 即可.
\end{proof}

\begin{proposition}(Weierstrass 定理的序列形式)\label{weierstrass序列形式}
	如果 $f$ 是 $[a,b]$ 上的一个连续函数,那么存在多项式 $P_n$ 的序列,使得
	\[
	\lim_{n\to\infty}P_n(x)=f(x)
	\]
	在 $[a,b]$ 上一致地成立.如果 $f$ 是实函数,那么 $P_n$ 可以取为实多项式.
\end{proposition}

\begin{proof}
	由定理\ref{weierstrass第一逼近定理},取 $\varepsilon=\dfrac1n$,则存在多项式 $P_n(x)$,使得对一切 $x\in[a,b]$,
	\[
	|P_n(x)-f(x)|<\frac1n.
	\]
	于是 $P_n\rightrightarrows f$ 于 $[a,b]$.
	
	反过来,若已经有多项式列 $P_n\rightrightarrows f$,则对任意 $\varepsilon>0$,存在 $N$,使得当 $n>N$ 时
	\[
	\sup_{x\in[a,b]}|P_n(x)-f(x)|<\varepsilon.
	\]
	取 $P=P_{N+1}$,便得到定理\ref{weierstrass第一逼近定理} 的形式.
\end{proof}

下面给出定理\ref{weierstrass序列形式} 的另一种证明,也就是 Rudin 中的证明方法.

\begin{proof}[另一种证明]
	仍不妨设 $[a,b]=[0,1]$.还可以假定 $f(0)=f(1)=0$. 因为若一般情形已经化为这种情形,只需考虑
	\[
	g(x)=f(x)-f(0)-x[f(1)-f(0)],\qquad 0\leqslant x\leqslant 1.
	\]
	则 $g(0)=g(1)=0$,而 $f-g$ 是一次函数,所以 $g$ 若为一致收敛多项式列的极限,则 $f$ 也是.
	
	在 $[0,1]$ 以外的点上定义 $f(x)=0$,则 $f$ 在整个实轴上一致连续.令
	\[
	Q_n(x)=c_n(1-x^2)^n\qquad (n=1,2,3,\cdots),
	\]
	其中 $c_n$ 由下式确定:
	\[
	\int_{-1}^{1}Q_n(x)\dd x=1.
	\]
	于是
	\[
	\frac1{c_n}=\int_{-1}^{1}(1-x^2)^n\dd x
	\geqslant \int_{-1/\sqrt n}^{1/\sqrt n}(1-x^2)^n\dd x
	\geqslant \int_{-1/\sqrt n}^{1/\sqrt n}(1-nx^2)\dd x
	=\frac{4}{3\sqrt n}.
	\]
	从而
	\[
	c_n<\sqrt n.
	\]
	上式中用到的不等式 $(1-x^2)^n\geqslant 1-nx^2$ 是显然的.这说明对任意 $\delta>0$,
	\[
	Q_n(x)\leqslant \sqrt n(1-\delta^2)^n,
	\qquad \delta\leqslant |x|\leqslant 1,
	\]
	于是 $Q_n(x)\to0$ 在 $\delta\leqslant |x|\leqslant1$ 上一致成立.
	
	现在令
	\[
	P_n(x)=\int_{-1}^{1}f(x+t)Q_n(t)\dd t,
	\qquad 0\leqslant x\leqslant1.
	\]
	利用 $f$ 的定义并作变量代换,可得
	\[
	\begin{aligned}
		P_n(x)&=\int_{x-1}^{x+1}f(u)Q_n(u-x)\dd u\\
		&=\int_0^1 f(u)Q_n(u-x)\dd u.
	\end{aligned}
	\]
	由于 $Q_n(u-x)=c_n[1-(u-x)^2]^n$ 关于 $x$ 是多项式,所以 $P_n(x)$ 也是关于 $x$ 的多项式.
	
	给定 $\varepsilon>0$.由 $f$ 的一致连续性,存在 $\delta>0$,使得当 $|t|<\delta$ 时,
	\[
	|f(x+t)-f(x)|<\frac{\varepsilon}{2}
	\]
	对所有 $x$ 成立.记 $M=\max |f|$.由 $\int_{-1}^{1}Q_n(t)\dd t=1$,有
	\[
	\begin{aligned}
		|P_n(x)-f(x)|
		&\leqslant \int_{-1}^{1}|f(x+t)-f(x)|Q_n(t)\dd t\\
		&\leqslant \int_{-\delta}^{\delta}|f(x+t)-f(x)|Q_n(t)\dd t
		+\int_{\delta\leqslant |t|\leqslant1}|f(x+t)-f(x)|Q_n(t)\dd t\\
		&\leqslant \frac{\varepsilon}{2}+2M\int_{\delta\leqslant |t|\leqslant1}Q_n(t)\dd t.
	\end{aligned}
	\]
	而后一项随 $n\to\infty$ 一致趋于 $0$,故对充分大的 $n$,
	\[
	|P_n(x)-f(x)|<\varepsilon
	\]
	对一切 $x\in[0,1]$ 成立.
\end{proof}